\chapter{Contexto del problema y estado del arte}
\label{cap:2}

\section{Gestion de emergencias y primera decision}

La gestion de emergencias comprende el conjunto de actividades orientadas a prevenir, preparar, responder y recuperar la normalidad frente a sucesos que pueden comprometer la vida, los bienes, el medio ambiente o servicios esenciales. A diferencia de otros procesos de decision publica, las emergencias se caracterizan por presion temporal, informacion incompleta, incertidumbre, multiplicidad de actores y necesidad de coordinacion inmediata \cite{Quarantelli1988,Comfort2007}.

En este contexto, la primera decision operativa tiene un peso especial. No cierra la gestion del incidente, pero condiciona la atencion inicial, la movilizacion de recursos, el seguimiento y la posible escalada. El operador debe interpretar rapidamente se\~nales parciales, distinguir incidentes leves de incidentes criticos y mantener una trazabilidad minima que permita justificar la decision adoptada.

El concepto de conciencia situacional ayuda a explicar esta necesidad. Endsley \citeyear{Endsley1995} la define como la percepcion de elementos del entorno, la comprension de su significado y la proyeccion de su evolucion. En un centro 112, esa conciencia situacional depende de transformar texto, categorias preliminares, coordenadas y contexto normativo en una representacion operativa util.

\section{Sistemas de apoyo a la decision en emergencias}

Los sistemas de apoyo a la decision (Decision Support Systems, DSS) buscan asistir a personas o equipos en decisiones complejas. En emergencias, pueden organizar informacion, visualizar incidentes, simular escenarios, priorizar recursos o facilitar coordinacion \cite{Turoff2004,Comes2011}. Sin embargo, muchos DSS se centran en cuadros de mando, mapas o gestion de recursos, sin abordar de forma especifica la priorizacion explicable de incidentes individuales en un servicio 112.

El prototipo de este TFM se situa en una franja mas acotada: apoyo a la primera decision mediante una recomendacion P1--P4, trazable y revisable. No pretende sustituir sistemas de despacho, mapas de riesgo o plataformas institucionales. Su aportacion consiste en conectar informacion inicial del incidente, variables operativas, modelo interpretable y explicacion normativa.

\section{Procesamiento del lenguaje natural en emergencias}

El procesamiento del lenguaje natural (NLP) se ha utilizado ampliamente para extraer informacion de mensajes de crisis, partes breves, redes sociales o comunicaciones ciudadanas. Trabajos clasicos de \textit{crisis informatics} mostraron que los textos no estructurados pueden contener localizaciones, peticiones de ayuda, da\~nos, necesidades y se\~nales de riesgo \cite{Vieweg2010,Caragea2011,Imran2015}.

Los modelos recientes basados en transformers han mejorado la representacion semantica de textos de crisis. CrisisTransformers, por ejemplo, propone modelos preentrenados y codificadores especializados para clasificacion, busqueda semantica y agrupacion de textos relacionados con crisis \cite{Lamsal2024}. Para este TFM, esa linea es relevante porque confirma la utilidad del NLP, pero tambien marca una frontera: procesar texto no equivale a decidir prioridad operativa.

En el sistema propuesto, el NLP es la Capa 1. Su funcion es extraer variables V01--V15 y se\~nales textuales como \texttt{signal\_herido\_grave}, \texttt{signal\_atrapado}, \texttt{signal\_intoxicacion}, \texttt{signal\_incendio} o \texttt{signal\_rescate}. La prioridad no la decide el extractor textual, sino la Capa 2 interpretable, lo que reduce el riesgo de convertir un clasificador de texto en una caja negra de decision. En la version v0.1.0 esta extraccion es determinista, pero la linea futura contempla modelos transformer en espa\~nol como MarIA/RoBERTa-bne \cite{GutierrezFandino2022MarIA} para reforzar la deteccion semantica de se\~nales.

\section{Supervision debil y construccion de etiquetas P1--P4}

Una dificultad central del proyecto es que el Registro de Emergencias del 112 de Castilla y Leon no contiene una prioridad oficial P1--P4. Esto impide entrenar directamente un modelo supervisado convencional sobre una verdad institucional. La solucion adoptada es construir una etiqueta academica mediante supervision debil.

La supervision debil permite combinar varias fuentes imperfectas de etiquetado, como reglas heuristicas, anotadores automaticos, clustering, modelos auxiliares o criterios expertos, para generar una etiqueta probabilistica de entrenamiento. Snorkel es una referencia consolidada en esta linea, al formalizar el uso de funciones de etiquetado y modelos generativos para crear datos de entrenamiento a partir de supervision imperfecta \cite{Ratner2020Snorkel}.

En este TFM, la supervision debil no se presenta como sustituto de una validacion experta institucional. Es un mecanismo academico y reproducible para poder entrenar y evaluar el prototipo con datos reales abiertos. Por ello se exige medir acuerdo, documentar la guia de etiquetado, analizar discrepancias y realizar ablaciones que eviten circularidad entre reglas de etiquetado y reglas del modelo.

\section{RuleFit, interpretabilidad y comparacion con modelos opacos}

La Capa 2 del sistema utiliza RuleFit como nucleo interpretable. RuleFit combina reglas derivadas de arboles con terminos lineales y seleccion regularizada, produciendo un conjunto de reglas con pesos que puede inspeccionarse y limitarse \cite{Friedman2008RuleFit}. Esta propiedad es especialmente adecuada para emergencias: el operador y el tribunal no solo necesitan saber que el sistema recomienda P1, sino que deben poder ver que reglas se han activado y por que.

El proyecto mantiene tambien un baseline experto y una ejecucion diagnostica con \texttt{imodels}. El baseline experto permite disponer de una referencia transparente, construida a partir de conocimiento normativo y se\~nales operativas. La comparativa con \texttt{imodels} permite contrastar la familia RuleFit recomendada con una implementacion externa, sin sustituir el nucleo reproducible seleccionado para v0.1.0. Esta separacion evita presentar un modelo opaco o poco controlado como decision operativa en un dominio de alto impacto.

La inclusion de un baseline experto basado en reglas heuristicas tradicionales y de un modelo alternativo tipo One-vs-Rest mediante \texttt{imodels.RuleFitClassifier} permite delimitar la ventaja analitica de la propuesta. El contraste experimental valida que la version optimizada RuleFit-lite seleccionada no solo mejora el rendimiento global, sino que reduce significativamente el coste computacional frente a implementaciones canonicas de la literatura.

La literatura sobre interpretabilidad refuerza esta decision. Rudin \citeyear{Rudin2019} defiende que, en decisiones de alto riesgo, debe preferirse un modelo interpretable desde el dise\~no cuando sea viable, en lugar de explicar a posteriori una caja negra. Esta idea es coherente con la arquitectura del TFM: RuleFit decide la prioridad; el LLM explica y comunica, pero no decide de forma autonoma.

\section{Explicabilidad, LLM local, RAG y MCP}

La explicabilidad en este TFM no se reduce a mostrar una probabilidad. La recomendacion debe acompa\~narse de reglas activadas, confianza, advertencias cuando proceda y citas normativas. Para ello, la Capa 3 utiliza un LLM local cuantizado \cite{Llama3_2024,Qwen25_2024}, recuperacion aumentada sobre un corpus normativo de Castilla y Leon y herramientas MCP.

\begin{figure}[htbp]
\centering
\fbox{%
\begin{minipage}{0.92\textwidth}
\centering
\small
\textbf{Entrada de texto del incidente}
$\rightarrow$
\textbf{Extraccion de se\~nales (Capa 1)}
$\rightarrow$
\textbf{Recomendacion y reglas activas (Capa 2)}
$\rightarrow$
\textbf{Orquestador RAG + prompt (Capa 3)}
$\rightarrow$
\textbf{Explicacion para el operador (HITL)}
\end{minipage}
}
\caption{Flujo conceptual XAI: separacion entre modelo decisor interpretable y capa generativa explicativa.}
\label{fig:flujo-xai-rulefit-rag}
\end{figure}

El uso de un LLM local responde a dos necesidades: generar explicaciones comprensibles para el operador y preservar soberania del dato. La recuperacion aumentada (RAG) \citeA{Lewis2020RAG} evita que la explicacion dependa solo de memoria generativa: el sistema consulta fragmentos del corpus normativo, recuperados mediante embeddings tipo Sentence-BERT \cite{Reimers2019SBERT}, y los utiliza para sustentar citas legales. En el prototipo real se utiliza el codificador de sentencias \texttt{paraphrase-multilingual-MiniLM-L12-v2}, por su equilibrio entre soporte multilingue, cobertura para textos en espa\~nol e ingles y ligereza operativa en ejecucion local. MCP aporta un mecanismo estandar para exponer herramientas como busqueda normativa, detalle de reglas y cita legal de prioridades; \citeA{Anthropic2024MCP} lo presento como un protocolo abierto para conectar asistentes de IA con fuentes de datos y herramientas externas.

El uso de un codificador de sentencias ligero como \texttt{paraphrase-multilingual-MiniLM-L12-v2}, con 384 dimensiones de \textit{embedding}, permite resolver busquedas semanticas sobre el corpus normativo local con baja latencia, mantener el consumo de memoria dentro del perfil de hardware local y asegurar que las explicaciones normativas se nutran de fragmentos legales exactos de forma reproducible.

Esta capa debe entenderse con precision metodologica. El LLM no decide la prioridad, no la recalcula y no debe inventar fundamento juridico. Su funcion es traducir una recomendacion ya calculada por Capa 2 en una explicacion breve, trazable y revisable. La fundamentacion juridica se obtiene de fragmentos recuperados del corpus normativo y de las reglas activadas por el modelo interpretable, de modo que la orquestacion no convierte al LLM en una caja negra decisoria.

El riesgo de alucinacion se mitiga mediante tres salvaguardas implementadas en el prototipo. Primero, la temperatura de inferencia se fija estrictamente en \texttt{0.0}. Segundo, la busqueda RAG queda acotada a fragmentos del corpus BOCYL/BOE incorporado al repositorio. Tercero, si el modelo local no responde o supera el SLA, se activa un modo degradado determinista que genera explicaciones automaticas a partir de plantillas, reglas activadas y metadatos de Capa 2.

\section{Sistemas y referencias relacionadas}

Existen varias familias de trabajos y sistemas relacionados. AIDR demostro la utilidad de combinar aprendizaje automatico y participacion humana para clasificar mensajes de crisis en redes sociales \cite{Imran2014AIDR}. TREC Incident Streams proporciono benchmarks para clasificar mensajes de crisis y estimar criticidad \cite{TRECIS2021}. GDACS integra alertas globales multirriesgo e informacion para coordinacion internacional \citeyear{GDACS2025}. Copernicus Emergency Management Service aporta cartografia rapida y productos geoespaciales para apoyar la respuesta \cite{CopernicusEMS2025}. Las revisiones recientes del JRC y SAPEA subrayan el potencial de la IA en gestion de emergencias, pero tambien la necesidad de control humano, interpretabilidad y robustez \cite{JRC2025AI,SAPEA2025}.

Ninguna de estas referencias resuelve exactamente el problema abordado aqui: priorizar incidentes 112 individuales en Castilla y Leon mediante una escala academica P1--P4, con variables trazables, RuleFit interpretable, explicacion normativa y supervision humana.

\begin{table}[htbp]
\centering
\scriptsize
\setlength{\tabcolsep}{3pt}
\begin{tabular}{|p{2.6cm}|p{3.3cm}|p{3.5cm}|p{3.5cm}|}
\hline
\textbf{Sistema o enfoque} & \textbf{Que aporta} & \textbf{Limitacion frente al TFM} & \textbf{Relacion con la propuesta} \\ \hline
AIDR & Clasificacion colaborativa de mensajes de crisis & Mensajes de redes sociales, no incidentes 112 consolidados & Inspira colaboracion humano-maquina \\ \hline
TREC-IS & Criticidad e informacion util en crisis & Prioriza mensajes, no incidentes con normativa local & Justifica tareas de criticidad y evaluacion \\ \hline
CrisisTransformers & NLP especializado para textos de crisis & No incorpora prioridad P1--P4 ni reglas normativas & Referente para Capa 1 NLP \\ \hline
GDACS & Alerta global multirriesgo & Escala internacional, no sala 112 autonoma & Inspira amenaza, exposicion y vulnerabilidad \\ \hline
Copernicus EMS & Cartografia y productos geoespaciales & No decide prioridad ni explica normativa & Fuente futura de contexto geoespacial \\ \hline
RuleFit & Modelo interpretable basado en reglas ponderadas & Requiere etiquetas y control de sparsity & Nucleo de Capa 2 \\ \hline
LLM + RAG + MCP & Explicacion y acceso a corpus/herramientas & Riesgo de alucinacion si no se acota & Capa 3 explicativa, no decisoria \\ \hline
\end{tabular}
\caption{Comparacion sintetica entre enfoques relacionados y propuesta del TFM.}
\label{tab:comparacion-sistemas-existentes}
\end{table}

\section{Brecha identificada}

A partir de la revision anterior, la brecha del TFM puede formularse asi: falta un prototipo academico, reproducible y explicable que combine incidentes reales 112, etiqueta P1--P4 academica, NLP en espa\~nol, supervision debil, RuleFit interpretable, explicacion normativa con LLM local, herramientas MCP y evaluacion orientada a seguridad del operador.

La propuesta cubre esa brecha con una arquitectura sobria. La Capa 1 estructura el incidente; la Capa 2 recomienda prioridad; la Capa 3 explica; el operador decide; y el sistema registra todo para auditoria. Esta cadena preserva el control humano y evita que la IA aparezca como autoridad autonoma.

\section{Conclusiones del estado del arte}

La revision permite extraer cinco conclusiones. Primera, la gestion de emergencias exige sistemas que ayuden a priorizar bajo incertidumbre, no solo a acumular informacion. Segunda, el NLP es necesario para estructurar texto, pero no suficiente para decidir prioridad. Tercera, la ausencia de etiquetas oficiales P1--P4 obliga a construir una etiqueta academica mediante supervision debil y validacion explicita. Cuarta, RuleFit ofrece un equilibrio adecuado entre aprendizaje y explicabilidad para el nucleo de decision. Quinta, los LLM locales con RAG y MCP pueden mejorar la comunicacion de la recomendacion, siempre que se mantengan subordinados al modelo interpretable y a la supervision humana.

Estas conclusiones preparan el transito al Capitulo 3, donde se desarrolla el marco operativo y normativo de Castilla y Leon que fundamenta las variables, reglas y citas del sistema.
